\section{导读：为什么核聚变又回到 AI 时代的中心}

本节先建立整期的学习目标。这期访谈表面上讲的是核聚变创业，实际牵出一个更大的问题：当 AI、数据中心和未来 AGI 持续扩大能源需求时，文明的上限可能不再只由算法、芯片和资本决定，还会由稳定、便宜、可扩展的能源供给决定。张小珺用一句很有记忆点的话开场：AGI 或许不怕人类，但一定怕断电。

杨钊给出的视角适合做一份科普型技术笔记。他不是只讲愿景，而是把可控核聚变拆成基础物理、托卡马克路线、三重积和 Q 值、高温超导材料、装置成本、创业组织形态、资金消耗、AI 对聚变的作用，以及聚变电站真正商业化还要跨过的工程台阶。读这期时，不要只问“什么时候实现无限能源”，更要问“哪些物理指标、工程系统和商业成本共同决定它能不能成为电力产品”。

\begin{importantbox}{本期核心命题}
可控核聚变的科学可行性已经有长期积累，但商业化的关键是把度电成本降到接近甚至低于火电。高温超导托卡马克的意义，是用更高磁场缩小装置尺寸，从而把一个大科学装置问题压缩成可被创业公司推进的工程和成本问题。
\end{importantbox}

\begin{knowledgebox}{视觉策略说明}
本视频是固定访谈画面，没有 slides、白板或产品演示。正文只使用封面作为来源识别，正文图像全部为自制概念图，用来解释能量链、托卡马克系统栈、聚变瓶颈、AI 与能源约束、创业路线图和投资逻辑。
\end{knowledgebox}

\begin{figure}[H]
\centering
\includegraphics[width=0.96\textwidth]{energy-ai-loop.png}
\caption{AI 与能源约束：AGI 不一定怕人类，但一定怕断电。自制概念图，依据 00:01:46--00:02:43 与 01:44:42--01:50:00 对谈内容整理。}
\end{figure}

\begin{knowledgebox}{读图：AI 和核聚变不是简单跨界热点}
图中从 AI 到 Data Center，再到 Energy、Fusion 和 Civilization，表达的是一个资源闭环：模型越强，电力需求越大；能源成本越低，计算、制造、交通和太空活动的上限越高。核聚变对 AI 的意义不只是“供电”，而是可能改变长期算力增长的边界条件。
\end{knowledgebox}

\subsection{本章小结}

这期的主线是：核聚变不是一个遥远科幻口号，而是一个由物理指标、工程可行性、装置成本和资本耐心共同决定的能源产品化问题。

